\documentclass[11pt,a4paper]{article}
\usepackage[margin=22mm]{geometry}
\usepackage{amsmath,amssymb,booktabs}
\usepackage[hidelinks]{hyperref}
\setlength{\parindent}{0pt}
\setlength{\parskip}{6pt}
\title{How the Combined Bus Ripple Is Calculated}
\author{Five-channel boost-converter calculator}
\date{14 October 2026}

\begin{document}
\maketitle

\section{The idea in plain English}

The load wants a steady current. Each boost converter supplies the output in
pulses. Whenever the converters supply less than the load wants, the output
capacitor makes up the difference and its voltage falls. Whenever they supply
more, the extra current charges the capacitor and its voltage rises.

The calculation therefore does three things:
\begin{enumerate}
  \item Work out when each converter supplies current to the output.
  \item Add those currents and subtract the load current.
  \item Accumulate the resulting charge and divide by capacitance to get voltage.
\end{enumerate}

\textbf{Interleaving changes the timing of the pulses, not their average power.}
The comparison holds the operating points and capacitance fixed. The current
calculator predicts
\[
  \boxed{\Delta V_{\mathrm{pp,aligned}}=101.07\ \mathrm{mV}},\qquad
  \boxed{\Delta V_{\mathrm{pp,72^\circ}}=12.11\ \mathrm{mV}}.
\]
These are predictions of the simplified model below, not measured board results.
``Peak-to-peak'' means the largest voltage minus the smallest voltage over a
complete switching period. It does not mean the peak excursion from the mean.

\section{Inputs and assumptions}

The values in \texttt{calculator.py} are
\begin{align*}
  V_{\mathrm{bus}} &=48\ \mathrm{V}, &
  f_s &=300\ \mathrm{kHz}, & T_s=1/f_s&=3.333333\ \mu\mathrm{s},\\
  \eta &=0.95, & L&=33\ \mu\mathrm{H}, & G&=1000\ \mathrm{W/m^2}.
\end{align*}
Each converter is assigned $48.84\ \mu\mathrm{F}$ of output capacitance.
For the combined calculation, we explicitly assume that all five capacitors
behave as one directly parallel bank:
\[
 C_{\mathrm{bus}}=5\times48.84\ \mu\mathrm{F}=244.2\ \mu\mathrm{F}.
\]
This is an editable modelling assumption. Output isolation, wiring and actual
capacitor placement must be considered before treating it as a board fact.

The model assumes continuous-conduction boost operation, constant input and bus
voltages within a switching period, and constant load current. It uses ideal
boost duty ratios and applies efficiency to average output power. The combined
calculation approximates output pulses as rectangular: inductor ripple,
capacitor ESR/ESL, dead time, output isolation and battery impedance are omitted.
The 48 V setting establishes the operating point; it is \emph{not} an ideal
voltage source clamping the simulated ripple to zero.

\section{From the panel to each converter's operating point}

Let $k$ identify F1, BL, BR, R1 or R2. The panel model keeps its MPP voltage
fixed and scales its reference MPP current linearly with irradiance:
\[
 V_{\mathrm{in},k}=V_{\mathrm{mpp},k},\qquad
 I_{\mathrm{in},k}=I_{\mathrm{mpp,ref},k}
       \frac{G_k}{1000\ \mathrm{W/m^2}}.
\]
All five panels currently have $I_{\mathrm{mpp,ref}}=8.22\ \mathrm{A}$, so
at the selected irradiance all five input currents equal $8.22\ \mathrm{A}$.
The stored open-circuit voltage and short-circuit current are not used in this
particular operating-point calculation.

The ideal boost voltage relation gives the low-side switch's on-time fraction:
\[
 V_{\mathrm{bus}}=\frac{V_{\mathrm{in},k}}{1-D_k}
 \quad\Longrightarrow\quad
 \boxed{D_k=1-\frac{V_{\mathrm{in},k}}{V_{\mathrm{bus}}}}.
\]
For example, F1 uses $D=1-27.59/48=0.42520833$, or $42.520833\%$.
Efficiency is not inserted into this ideal duty relation in the current code.

Average power and output current are
\[
 P_{\mathrm{in},k}=V_{\mathrm{in},k}I_{\mathrm{in},k},\qquad
 P_{\mathrm{out},k}=\eta P_{\mathrm{in},k},\qquad
 \boxed{\overline I_{\mathrm{out},k}
       =\frac{\eta V_{\mathrm{in},k}I_{\mathrm{in},k}}{V_{\mathrm{bus}}}}.
\]

\begin{center}
\begin{tabular}{lrrrr}
\toprule
Panel & $V_{\mathrm{in}}$ (V) & $I_{\mathrm{in}}$ (A) & $D$ & $\overline I_{\mathrm{out}}$ (A)\\
\midrule
F1 & 27.59 & 8.22 & 0.42520833 & 4.488548125\\
BL & 16.15 & 8.22 & 0.66354167 & 2.627403125\\
BR & 16.15 & 8.22 & 0.66354167 & 2.627403125\\
R1 & 21.54 & 8.22 & 0.55125000 & 3.504288750\\
R2 & 12.11 & 8.22 & 0.74770833 & 1.970145625\\
\bottomrule
\end{tabular}
\end{center}

The constant load is chosen to consume the total average output current:
\[
 \boxed{I_{\mathrm{load}}=\sum_k\overline I_{\mathrm{out},k}
                         =15.21778875\ \mathrm{A}}.
\]
This makes average capacitor current zero, so the bus does not drift upward or
downward from cycle to cycle.

\section{The individual ripple and peak values printed by the calculator}

During switch-on the inductor sees approximately $V_{\mathrm{in}}$.
Using $v_L=L\,di_L/dt$ for an on-time of $D_kT_s$ gives
\[
 \boxed{\Delta I_{L,k,\mathrm{pp}}
       =\frac{V_{\mathrm{in},k}D_k}{L f_s}},\qquad
 \boxed{I_{L,k,\mathrm{peak}}
       =I_{\mathrm{in},k}+\frac{\Delta I_{L,k,\mathrm{pp}}}{2}}.
\]
The average inductor current equals input current in this boost approximation.
For F1 these equations give $1.1850\ \mathrm{A}$ peak-to-peak and
$8.8125\ \mathrm{A}$ peak. The CCM assumption requires the valley current
$I_{\mathrm{in},k}-\Delta I_{L,k,\mathrm{pp}}/2$ to remain positive.

For one converter considered alone, its capacitor supplies the load for
$D_kT_s$ while the switch is on. The charge removed is
$\Delta Q=\overline I_{\mathrm{out},k}D_kT_s$, hence
\[
 \boxed{\Delta V_{\mathrm{out},k,\mathrm{pp}}
       \simeq\frac{\overline I_{\mathrm{out},k}D_k}{f_s C_{\mathrm{out},k}}}.
\]
For F1, using its own $48.84\ \mu\mathrm{F}$ gives $130.26\ \mathrm{mV}$.
This is the familiar capacitive-ripple approximation, without ESR.
\textbf{The combined routine does not add these individual ripple voltages.}
It adds the output-current waveforms and integrates into the shared capacitance.
Also, the individual inductor-ripple values printed above are currently
\emph{not fed into} the combined bus-ripple model.

\section{Build the output-current pulses}

When the low-side switch is on, the boost inductor is being charged and its
output path supplies no current in this model. When the switch turns off,
the converter supplies the output for the remaining fraction $1-D_k$.

Replace this output interval with a constant pulse of amplitude $A_k$.
Its cycle average must equal the output current already calculated:
\[
 \overline I_{\mathrm{out},k}=A_k(1-D_k)
 \quad\Longrightarrow\quad
 \boxed{A_k=\frac{\overline I_{\mathrm{out},k}}{1-D_k}}.
\]
For these inputs, substitution gives a useful simplification:
\[
 A_k=\frac{\eta V_{\mathrm{in},k}I_{\mathrm{in},k}/V_{\mathrm{bus}}}
               {V_{\mathrm{in},k}/V_{\mathrm{bus}}}
     =\eta I_{\mathrm{in},k}=0.95\times8.22=7.809\ \mathrm{A}.
\]
Thus all five rectangular pulses have the same height, although their widths
and average output powers differ. This equality is particular to the common
input current and efficiency selected here.

\textbf{This is an effective, power-consistent pulse, not a detailed physical
inductor waveform.} A real CCM output interval follows the sloping inductor
current. The use of ideal duty together with efficiency-scaled output power
makes the effective height $7.809\ \mathrm{A}$ rather than $8.22\ \mathrm{A}$.
A more detailed loss and switching model would need to reconcile both voltage
and current waveforms physically.

\section{Put each pulse at the requested phase}

Use normalized time $x=t/T_s$, so one switching cycle runs from $0$ to $1$.
Convert the phase angle into a fraction of a cycle:
\[
 p_k=\frac{\phi_k\bmod360^\circ}{360^\circ},\qquad
 u_k(x)=(x-p_k)\bmod1.
\]
Here ``modulo 1'' means wrap back to the beginning after one complete cycle.
Positive phase delays the low-side turn-on. The output pulse is
\[
 i_{\mathrm{out},k}(x)=
 \begin{cases}
  0, & 0\leq u_k(x)<D_k,\\
  A_k, & D_k\leq u_k(x)<1.
 \end{cases}
\]
Its start is $s_k=(p_k+D_k)\bmod1$ and its end is $p_k$ in the next local cycle.
If the interval crosses the global cycle boundary, its current is already
present at $x=0$. The code accounts for that; it does not discard the wrapped
part of a pulse. At exactly zero duty the output is continuous.

For the two cases, in panel order F1, BL, BR, R1, R2:
\begin{align*}
 \text{Aligned:}\quad &\boldsymbol\phi=(0,0,0,0,0)^\circ,\\
 \text{Interleaved:}\quad &\boldsymbol\phi=(0,72,144,216,288)^\circ,\\
 &\boldsymbol p=(0,0.2,0.4,0.6,0.8).
\end{align*}
Adjacent starts are $T_s/5=0.666667\ \mu\mathrm{s}$ apart.
In the interleaved case the output-pulse starts are
\[
 \boldsymbol s=(0.42520833,\ 0.86354167,\ 0.06354167,\ 0.15125,\ 0.54770833).
\]

\section{Convert the current imbalance into bus voltage}

Kirchhoff's current law at the bus gives
\[
 \boxed{i_C(x)=\sum_k i_{\mathrm{out},k}(x)-I_{\mathrm{load}}}.
\]
Positive $i_C$ charges the capacitor; negative $i_C$ discharges it. Since
$i_C=C_{\mathrm{bus}}\,dv/dt$,
\[
 v(x)-v(0)=\frac{1}{f_s C_{\mathrm{bus}}}\int_0^x i_C(\xi)\,d\xi.
\]
The pulses only change at switching edges. Sort all those edges and integrate
each constant-current interval exactly:
\[
 \boxed{v_{j+1}=v_j+
        \frac{i_{C,j}(x_{j+1}-x_j)}{f_s C_{\mathrm{bus}}}}.
\]
There is no time-step approximation in this event-based integration of the
rectangular-pulse model. Voltage is linear between edges, so its extrema occur
at the interval boundaries. Finally,
\[
 \boxed{\Delta V_{\mathrm{bus,pp}}=\max_j(v_j)-\min_j(v_j)}.
\]
We set $v(0)=0$ merely as a voltage reference. These numbers are changes from
the start of the cycle, not absolute bus voltage or necessarily changes from
the mean. A vertical offset does not affect peak-to-peak ripple.

Here $f_s C_{\mathrm{bus}}=300000\times244.2\times10^{-6}=73.26$.
For convenient table arithmetic,
\[
 \Delta v_j\ [\mathrm{mV}]
 =\frac{1000}{73.26}\,i_{C,j}\ [\mathrm{A}]\,(x_{j+1}-x_j).
\]

\section{Worked case: all five phases aligned}

At the start of the cycle every switch is on, so no converter supplies the
output. The capacitor initially supplies the entire $15.21778875\ \mathrm{A}$.
F1 starts supplying current first, followed by R1, then BL and BR together,
then R2. All arithmetic below uses unrounded input values.

\begin{center}
\small
\begin{tabular}{rrlrr}
\toprule
$x_j$ & $x_{j+1}$ & Supplying output & $i_C$ (A) & $v_{j+1}$ (mV)\\
\midrule
0.000000 & 0.425208 & None & -15.217789 & -88.325561\\
0.425208 & 0.551250 & F1 & -7.408789 & -101.072163\\
0.551250 & 0.663542 & F1, R1 & 0.400211 & -100.458726\\
0.663542 & 0.747708 & F1, BL, BR, R1 & 16.018211 & -82.055785\\
0.747708 & 1.000000 & All five & 23.827211 & 0.000000\\
\bottomrule
\end{tabular}
\end{center}

For example, the first interval gives
\[
 \Delta v_0=\frac{-15.21778875\times0.4252083333}{73.26}
           =-0.088325561\ \mathrm{V}.
\]
The second interval removes another $12.746602\ \mathrm{mV}$. After R1 begins
supplying current, the two active pulses supply $2\times7.809=15.618\ \mathrm{A}$,
slightly above the load, and the capacitor starts charging again. Therefore
\[
 \Delta V_{\mathrm{pp}}=0-(-101.072163)=101.072163\ \mathrm{mV}.
\]

\section{Worked case: 72-degree spacing}

At the global start of this cycle, BL is already in the output interval that
began in the previous cycle. Throughout this example, between one and three
converters supply the output. With $n$ active pulses,
\[
 i_C=n(7.809)-15.21778875\ \mathrm{A}.
\]
One pulse gives $-7.40878875\ \mathrm{A}$, two give
$+0.40021125\ \mathrm{A}$, and three give $+8.20921125\ \mathrm{A}$.
The capacitor spends much of the cycle very close to current balance.

\begin{center}
\small
\begin{tabular}{rrlrr}
\toprule
$x_j$ & $x_{j+1}$ & Supplying output & $i_C$ (A) & $v_{j+1}$ (mV)\\
\midrule
0.000000 & 0.063542 & BL & -7.408789 & -6.425973\\
0.063542 & 0.151250 & BL, BR & 0.400211 & -5.946832\\
0.151250 & 0.200000 & BL, BR, R1 & 8.209211 & -0.484110\\
0.200000 & 0.400000 & BR, R1 & 0.400211 & 0.608468\\
0.400000 & 0.425208 & R1 & -7.408789 & -1.940852\\
0.425208 & 0.547708 & F1, R1 & 0.400211 & -1.271648\\
0.547708 & 0.600000 & F1, R1, R2 & 8.209211 & 4.587939\\
0.600000 & 0.800000 & F1, R2 & 0.400211 & 5.680516\\
0.800000 & 0.863542 & F1 & -7.408789 & -0.745457\\
0.863542 & 1.000000 & F1, BL & 0.400211 & 0.000000\\
\bottomrule
\end{tabular}
\end{center}

The minimum is $-6.425973\ \mathrm{mV}$ and the maximum is
$+5.680516\ \mathrm{mV}$. Thus
\[
 \Delta V_{\mathrm{pp}}=5.680516-(-6.425973)=12.106489\ \mathrm{mV}.
\]
The reduction, using unrounded results, is
\[
 100\left(1-\frac{12.106489}{101.072163}\right)\%=88.02\%.
\]
The large reduction comes from shrinking the accumulated charge imbalance.
It is not an assumption that five phases automatically divide ripple by five.

\section{Why the combined ripple can be below one converter's ripple}

The individual F1 result uses $48.84\ \mu\mathrm{F}$; the combined result uses
$244.2\ \mu\mathrm{F}$. Those are different capacitor sizes. For $N$ identical,
aligned converters, each contributing current $\overline I$ and capacitance
$C$, the simple result is
\[
 \Delta V_{\mathrm{pp}}=\frac{N\overline I D}{f_s NC}
                      =\frac{\overline I D}{f_s C}.
\]
Both current and capacitance increase together, so ripple need not increase.
Our five channels also have different duties, so even aligned starts do not
mean aligned output-pulse starts.

For a fair phase comparison we keep $C_{\mathrm{bus}}$ identical in both cases.
If the effective total were instead only $48.84\ \mu\mathrm{F}$, both combined
results would be five times larger: approximately $505.36\ \mathrm{mV}$ aligned
and $60.53\ \mathrm{mV}$ interleaved. The percentage reduction would be unchanged
within this model.

\section{What to trust, and what this calculation does not establish}

The integration is exact for the stated rectangular-current model. The circuit
model itself is approximate. In particular:
\begin{itemize}
 \item Actual inductor-current slopes can matter substantially when the
 remaining combined ripple is small. The current model's bus result does not
 depend on $L$, even though $L$ affects the individual inductor-ripple display.
 \item ESR, ESL and switching edges introduce voltage ripple and spikes not
 represented by charge integration into an ideal capacitor.
 \item Battery impedance and capacitor placement determine where switching
 current flows. An ideal voltage source directly across the bus would clamp
 voltage ripple; a real battery and its wiring require an impedance model.
 \item The assumed parallel capacitance and its effective value under voltage
 bias must be checked against the hardware.
 \item Low irradiance can invalidate CCM. This model does not calculate DCM
 switching or reverse-current behaviour.
 \item Equal spacing is the requested comparison, not an optimisation result.
 With unequal duties, changing which panel receives each phase can change
 ripple. No claim of an optimal phase assignment is made here.
\end{itemize}

Consequently, $88.02\%$ is a model prediction at these operating points, not a
promised reduction on the real board.

\section{How this maps to the code}

\begin{description}
 \item[\texttt{solar\_panels.py}] Scales input MPP current with irradiance.
 \item[\texttt{converter.py}] Calculates duty, mean output current, inductor
 ripple, peak current and the single-converter voltage-ripple estimate.
 \item[\texttt{bus\_ripple.py}] Builds pulse start/end events, handles pulses
 wrapping across the cycle boundary, integrates capacitor current between
 sorted events and returns maximum minus minimum voltage.
 \item[\texttt{calculator.py}] Supplies the same operating points, frequency
 and capacitance to the zero-phase and 72-degree phase cases, then prints
 their peak-to-peak ripple and percentage reduction.
\end{description}

Checks performed include the single-converter analytical result, cancellation
for two equal 50\%-duty channels separated by $180^\circ$, cancellation for
five equal 80\%-duty channels separated by $72^\circ$, invariance to a common
phase shift, and comparison with independently sampled pulse waveforms.
These check the mathematics and implementation, not the physical assumptions.

\paragraph{Background reference.}
Texas Instruments, \emph{Basic Calculation of a Boost Converter's Power Stage},
SLVA372D, gives the single-converter inductor and output-capacitance relations:
\url{https://www.ti.com/lit/an/slva372d/slva372d.pdf}.
The event tables and combined-bus results above are derived from this project's
stated pulse model, not quoted from that application note.

\end{document}
